% Le due forze degli elastici sul dente (problema ud1-vett-ortodonzia): F_1 = 1,5 N
% lungo x, F_2 = 1,0 N a 60 gradi, e la somma R = (2,0; 0,87) N.
% Scala: 1 N = 2,4 cm (F_1 = 3,6, F_2 = 2,4, R = (4,8; 2,08)).
\begin{tikzpicture}[>=Stealth, font=\small]
  \draw[->] (-0.4,0) -- (5.6,0) node[right] {$x$};
  \draw[->] (0,-0.4) -- (0,2.8) node[above] {$y$};
  \coordinate (A) at (3.6,0);
  \coordinate (B) at (60:2.4);
  \coordinate (R) at (4.8,2.08);
  \draw[dashed, gray] (A) -- (R) -- (B);
  \draw[dashed, gray] (R) -- (R |- 0,0);
  \draw[->, very thick, formula] (0,0) -- (A);
  \node[formula, below] at (1.8,0) {$\vec F_1$};
  \draw[->, very thick, medica] (0,0) -- (B);
  \node[medica, left] at ($(0,0)!0.6!(B)$) {$\vec F_2$};
  \draw[->, very thick, esempio] (0,0) -- (R);
  \path (0,0) -- (R) node[esempio, pos=0.5, below, sloped] {$\vec R = \vec F_1 + \vec F_2$};
  \draw[thick, richiamo] (0.7,0) arc[start angle=0, end angle=60, radius=0.7];
  \node[richiamo] at (30:1.0) {\ang{60}};
  \draw[<->] (0,-0.75) -- (3.6,-0.75);
  \node[below] at (1.8,-0.75) {\SI{1,5}{\newton}};
  \draw[<->] (3.6,-0.75) -- (4.8,-0.75);
  \node[below] at (4.2,-0.75) {\SI{0,50}{\newton}};
  \node[right] at (4.8,1.04) {$R_y \approx \SI{0,87}{\newton}$};
\end{tikzpicture}
